\chapter*{Úvod}
\addcontentsline{toc}{chapter}{Úvod}

Cílém této práce je implementovat nástroj, který pro vstupní SAT formuli vygeneruje klauzule pro rozbíjení symetrií pro zvolenou množinu permutací na proměnných, a dále zkoumat, jak různé výběry množin permutací ovlivňují celkový počet rozbitých symetrií, množství přidaných klauzulí a celkově efektivitu nástroje.

Problém splnitelnosti booleovských formulí (SAT) je jedním ze základních rozhodovacích problémů. Cílem je určit, zda existuje ohodnocení proměnných zadané booleovské formule takové, že výsledný výraz je pravdivý. Přestože se nástroje pro řešení SAT (SAT solvery) instancí výrazně posunuly kupředu, složité instance často obsahují redundance a struktury, které stále příliš rozšiřují prohledávací prostor. Jedním ze zdrojů takové redundance jsou právě symetrie.

Symetrie v SAT formulích lze chápat jako permutace proměnných, které zachovávají pravdivost celé formule. Jako hlavní příklad symetrií v této práci budeme uvažovat záměnu ekvivalentních objektů ve vstupní úloze. Pro lepší představu si jako jednoduchý příklad uvedeme problém holubníku. Zde je cílem určit, zda je možné rozmístit zadaný počet holubů do zadaného počtu holubníků tak, že v každém holubníku bude nejvýše jeden holub. Na jednotlivé holuby ani holubníky nemáme žádné speciální požadavky, takže například pro $5$ holubníků a $1$ holuba máme $5$ různých řešení, protože nezáleží na tom, který holubník vybereme. Přitom z jednoho řešení na druhé se můžeme dostat permutací označení jednotlivých holubníků, protože jednotlivé holubníky jsou ekvivalentní. V tomto smyslu jsou i samotná řešení ekvivalentní a SAT solveru by stačilo zkontrolovat jediné z nich, aby přišel na to, jestli jsou všechna pravdivá, nebo nepravdivá. Avšak právě podobné permutované varianty stejného řešení SAT solvery často zbytečně procházejí mnohokrát během jednoho řešení, čímž se dramaticky zvětšuje prohledávací prostor. V ideálním případě SAT solver z každé třídy ekvivalencí řešení zkontroluje vždy nejvýše jedno řešení. K tomuto stavu se chceme přidáváním klauzulí co nejvíce přiblížit a říkáme mu úplné rozbití symetrie.

Úplné rozbití všech symetrií je ale v realitě většinou příliš náročné. SAT solver sice v takovém případě poběží kratší dobu, ale samotné generování všech potřebných klauzulí pro rozbíjení symetrií zabere příliš mnoho času. Musíme tedy generovat určitou podmnožinu klauzulí pro rozbíjení symetrií, která ale stále rozbije dostatečné množství symetrií. Tomuto přístupu říkáme částečné rozbití symetrie. \nocite{*}

Jedním z nejčastěji používaných způsobů rozbíjení symetrií je metoda Lex-Leader. Zde pro každou třídu ekvivalencí řešení definujeme jedno jako kanonické řešení. Toho docílíme přidáním omezujících podmínek, které splňuje pouze kanonické řešení. Základem je nejdříve si zvolit uspořádání na proměnných. Z tohoto uspořádání odvodíme uspořádání na celých ohodnoceních stejně jako u lexikografického uspořádání. Máme tedy

$$x_1 \prec x_2 \prec \dots \prec x_n$$

a pro každé ohodnocení $x$ a každou permutaci $\pi$ přidáme omezující podmínku 

$$x \leq_{\text{lex}} \pi(x).$$

Vybrané kanonické řešení je vždy lexikograficky nejmenší a tím pádem jako jediné z dané třídy ekvivalencí bude splňovat všechny přidané podmínky. Tímto zřejmě dosáhneme úplného rozbití symetrií. Jak již ale bylo řečeno, v tomto případě bychom potřebovali přidat příliš mnoho omezujících podmínek. Z tohoto důvodu se budeme zaměřovat na částečné rozbití symetrií, kdy se omezující podmínky generují pouze pro vybranou podmnožinu permutací.